% !TEX root = ../main.tex
%%%%%%%%%%%%%%%%%%%%%%%%%%%%%%%%%%%%%%%%%%%%%%%%%%%%%%%%%%%%%%%%%%%%%%%
% Results II - stress tests and the wage channel
%%%%%%%%%%%%%%%%%%%%%%%%%%%%%%%%%%%%%%%%%%%%%%%%%%%%%%%%%%%%%%%%%%%%%%%

\section{Results II: stress tests and the wage channel}
\label{sec:stress}

Each subsection below states a hypothesis that a critic might raise, turns on
the corresponding mechanism, and reports what was measured. The unifying
finding is not any one of them but what they share: a wage-driven feedback
that erodes capital, discovered in the first test and unmoved by every
instrument subsequently aimed at it, including the fiscal one of
\cref{sec:government}, which is reported separately because it also produces
a result of its own.

\subsection{Wage flexibility does not overturn the
\texorpdfstring{$\sigma = 1$}{sigma = 1} wage-led result, and it widens the
profit-led band}
\label{sec:wagecurve}

\paragraph{The objection.} The baseline holds the wage fixed, which suppresses
an offsetting channel: unemployment should depress the wage, make labor
cheaper, and slow the substitution toward capital. Turn that channel on and
the wage-led result should weaken.

\paragraph{What was measured.} The opposite of the objection happens, though
not in the obvious direction (\cref{tab:wagecurve}, \cref{fig:wagecurve}). $\sigstar$ rises, so the $\sigma = 1$ wage-led
result survives; $\sigma = 1$ stays above the frontier at every $\eta$,
while the empirical band $[0.40, 0.60]$ is pushed deeper into the
profit-led region. Wage flexibility does not rescue a wage-led reading of the
empirically supported band; it enlarges the band in which retention raises
output.

\begin{table}[htbp]
\centering
\small
\caption{$\sigstar(\eta)$ on $Y$ at $c_0 = 1.0$, estimated on the support
viable at every $\eta$ so that the columns are comparable. At
$\eta = 0.10$ and $\eta = 0.15$ the slope $\mathrm{d}Y/\mathrm{d}\rho$ crosses
zero twice over the support (\texttt{n\_crossings} $=2$ in
\texttt{ces\_b07\_sigma\_star.csv}); the entry is the first crossing, and is a
single point only at the two smaller $\eta$, where the crossing is unique.}
\label{tab:wagecurve}
\begin{tabular}{lrlrr}
\toprule
$\eta$ & $\sigstar$ & 95\% CI & $P(\sigstar > 0.60)$ & mean $U$\\
\midrule
0.00 & 0.654 & $[0.616,\,0.691]$ & 99.8\% & 0.528\\
0.05 & 0.666 & $[0.634,\,0.692]$ & 99.9\% & 0.543\\
0.10 & 0.725 & $[0.697,\,0.745]$ & 100\% & 0.565\\
0.15 & 0.740 & $[0.682,\,0.793]$ & 100\% & 0.579\\
\bottomrule
\end{tabular}
\end{table}

$\sigstar$ rises with $\eta$, moving further above the empirical range,
and mean unemployment rises with it ($0.53 \to 0.58$). The reason is that the
counter-channel the objection invokes is not the only one the wage curve
activates: a falling wage also shrinks the wage bill and hence demand. This is
the Kaleckian paradox of costs, and in this model it dominates substitution.
Wage flexibility does not auto-correct unemployment here; it makes it worse.

\begin{figure}[htbp]
\centering
\includegraphics[width=0.82\textwidth]{ces_b07_sigma_star_eta.png}
\caption{$\sigstar(\eta)$: wage flexibility moves the frontier away
from the empirical range.}
\label{fig:wagecurve}
\end{figure}

\paragraph{The price convention brackets this result.} The result above holds the
numéraire fixed at one, which makes the nominal wage the real wage, and that is
a convention, not a measurement. A normal-cost probe relaxes it: with constant
markup and productivity the price is $P_t = w_t/\bar w$, so a nominal wage cut is
passed through to the price and the real wage $w_t/P$ stays at $\bar w$. Under this
complete, instantaneous price pass-through $\sigstar$ rises more steeply with
$\eta$, not less (\cref{tab:prices}): the two conventions' intervals are cleanly
disjoint at $\eta = 0.15$ and only marginally so at $\eta = 0.10$, separated there
by about $0.001$. The mechanism is a Pigou effect on real balances (lower
prices raise real wealth, which raises demand) and it is stabilizing, so a
reader who took complete pass-through for the truth would conclude that wage
flexibility does auto-correct. Two cautions keep this a bracket rather than
a verdict. The probe deliberately occupies the most stabilizing corner of the
space (constant markup and productivity, no price stickiness), so it is an
upper bound on how far pass-through can stabilize. And the off control is not
itself resolved at twelve seeds: its $\sigstar$ is non-monotone in $\eta$, the
endpoint falling below the $\eta = 0.10$ value, and it carries two crossings at
$\eta \ge 0.10$ where the frontier is not a single point. The honest reading is a
bracket: at zero pass-through the wage curve does not auto-correct unemployment,
at complete pass-through it does, and the model cannot say where between them the
data sit. \Cref{sec:frontierlocal} returns to this as one of three senses in which
the frontier is a local object.

\begin{table}[htbp]
\centering
\small
\caption{The price convention brackets the wage-flexibility result:
$\sigstar(\eta)$ on $Y$ at $c_0 = 1.0$ with the numéraire fixed (off) and
under complete instantaneous pass-through $P_t = w_t/\bar w$ (on),
estimated on the support viable at every $\eta$. Vintage \texttt{ces\_b21}
(twelve seeds; an eleven-node $\sigma$ grid), a different vintage from the
\texttt{ces\_b07} of \cref{tab:wagecurve} and not commensurable with it, the
off column is the probe's own control, not \cref{tab:wagecurve}. \#cross
counts the sign crossings of $\mathrm{d}Y/\mathrm{d}\rho$ over the support; the
off control develops a second at $\eta \ge 0.10$, where $\sigstar$ is its first
crossing.}
\label{tab:prices}
\begin{tabular}{llrlr}
\toprule
$\eta$ & pass-through & $\sigstar$ & 95\% CI & \#cross\\
\midrule  % generated by scripts/make_tab_prices.py -- do not hand-edit
$0.00$  & off\,$\equiv$\,on & 0.648 & $[0.591,\,0.709]$ & 1\\
$0.05$  & off               & 0.680 & $[0.657,\,0.711]$ & 1\\
        & on                & 0.716 & $[0.671,\,0.770]$ & 1\\
$0.10$  & off               & 0.726 & $[0.700,\,0.750]$ & 2\\
        & on                & 0.806 & $[0.751,\,0.838]$ & 1\\
$0.15$  & off               & 0.711 & $[0.658,\,0.800]$ & 2\\
        & on                & 0.871 & $[0.843,\,0.909]$ & 1\\
\bottomrule
\end{tabular}
\end{table}

\paragraph{A collapse mechanism, discovered here.}
In the secondary regime $c_0 = 2.0$, wage flexibility destabilizes the
economy in the high-$\sigma$, low-$\rho$ corner: at $\sigma = 1.25$, 43\% of
seeds collapse to $Y \to 0$, $U \to 1$ at $\eta = 0.15$, and $\sigstar$ becomes
erratic or undefined. Tracing a collapsing trajectory identifies the
mechanism, which turns out to govern much of the rest of this paper:

(1) the wage oscillates, rising above $\bar w$ when $U \to 0$
(the $U_{\min}$ guard drives $w$ to $\approx 1.25$) and falling below it when
$U$ is high; (2) as $\sigma$ rises, $L^{\mathrm{profitmax}}$ becomes more
wage-sensitive, so the wage oscillation feeds a widening employment
oscillation; and (3) each cycle, investment fails to cover depreciation, so
capital erodes, and at low $\rho$ the economy eventually dies at $U = 1$.

In the empirically relevant region $\sigma \approx 0.5$ the same curve leaves
$w \approx \bar w$, produces no oscillation, and capital grows. The
collapse is capital erosion at high $\sigma$, not a monotone wage spiral, and
the subsistence floor $w_{\min}$ never stably binds, so this is genuine
viability failure, not a floor artefact. We name this the
wage\,$\to U \to$\,capital-erosion channel; \cref{sec:discussion}
returns to it.

\paragraph{The price probe separates the collapse.} Step~(2) of that mechanism
turns on the firm reacting to the real wage: at a fixed numéraire the
nominal wage is the real wage, so the wage oscillation is an oscillation in the
firm's own marginal-cost signal. Under complete pass-through it is not, 
$w_t/P \equiv \bar w$, and step~(2) does not start. Running $c_0 = 2.0$ with
prices on does not, however, abolish the collapse. At $\eta = 0$, where
$P \equiv 1$ by construction, the two conventions are identical (two cells
collapse, one of them fully, under each). Away from there the incidence that
climbed with $\eta$ under the fixed numéraire, cells with any collapse running
$2, 8, 12, 14$ across $\eta$, is flat with prices on, at $2, 3, 3, 2$;
excluding the degenerate $\eta = 0$ cell, cells with any collapse fall from $34$
to $8$ and the fully collapsed from $16$ to $3$. The probe therefore
separates the collapse into two components: a wage-driven part that
pass-through removes, and a residual, already present at $\eta = 0$ and unrelated
to wage flexibility, that it does not.

\subsection{Adaptive expectations move neither the frontier nor the collapse}
\label{sec:expectations}

\paragraph{The hypothesis.} A slower firm expectation should damp the
oscillation just identified, and might also select a different steady state.
We generalize the demand expectation from static to adaptive in the
partial-adjustment form of \citet{Nerlove1958},
$Y^e_t = Y^e_{t-1} + \lambda_e (D_{t-1} - Y^e_{t-1})$, which is the
constant-gain case of the adaptive-learning literature
\citep{EvansHonkapohja2001}. We sweep $\lambda_e$ rather than choosing it: that
literature supplies the functional form, not a level, and no reliable point
estimate exists for a firm's demand expectation in a model of this kind.

\paragraph{What was measured.} The frontier is $\lambda_e$-invariant within
confidence intervals (\cref{tab:expectations}): every point lies inside its
neighbours' intervals, and the empirical range stays below $\sigstar$ at every
gain, so it is profit-led at every gain, while $\sigma = 1$ stays above and
remains wage-led. What is robust to the expectation gain is the location
of the frontier, and with it both readings on either side of it; so is its rise
with $\eta$. There is no basin-selection effect.

\begin{table}[htbp]
\centering
\small
\caption{$\sigstar(\eta;\lambda_e)$ on $Y$, $c_0 = 1.0$, across-configuration
common support.}
\label{tab:expectations}
\begin{tabular}{llrlr}
\toprule
$\eta$ & $\lambda_e$ & $\sigstar$ & 95\% CI & $P(\sigstar > 0.60)$\\
\midrule
0.00 & 1.00 & 0.654 & $[0.616,\,0.691]$ & 99.8\%\\
0.00 & 0.50 & 0.686 & $[0.637,\,0.721]$ & 99.9\%\\
0.00 & 0.25 & 0.674 & $[0.639,\,0.709]$ & 99.9\%\\
0.10 & 1.00 & 0.725 & $[0.697,\,0.745]$ & 100\%\\
0.10 & 0.50 & 0.713 & $[0.667,\,0.752]$ & 100\%\\
0.10 & 0.25 & 0.721 & $[0.684,\,0.754]$ & 100\%\\
\bottomrule
\end{tabular}
\end{table}

The stabilization half of the hypothesis is not confirmed. The
collapse region is $\lambda_e$-invariant to within grid and seed noise (cells
with any collapse at $\eta = 0.10$: 16, 15, 14 for $\lambda_e = 1, 0.5, 0.25$,
with fully collapsed cells flat at 6; non-monotone at $\eta = 0.15$), and the
reference collapsing cell dies at $K = 0$, $U = 1$ at every gain. The
interpretation is now sharp: damping the demand expectation cannot
stabilize an instability that does not originate in demand.

\subsection{The accelerator's signal, smoothed}
\label{sec:invexpectations}

\paragraph{The hypothesis.} \Cref{sec:expectations} smoothed the firm's demand
expectation; the same question applies to the utilization signal the accelerator
itself reads. \Cref{eq:accelerator} reacts to realized utilization, and a
hypothesis fixed before the runs held that letting firms react instead to the
smoothed, adaptive $u^e_t = u^e_{t-1} + \lambda_u\,(u_{t-1} - u^e_{t-1})$ would
damp the accelerator and might move the sign result. We sweep the gain
$\lambda_u$ over $[0, 1]$; $\lambda_u = 1$ is the default of the main results,
where $u^e_t = u_{t-1}$ exactly.

\paragraph{What was measured.} The hypothesis is falsified on two levels. The
mechanism fails at its root: the temporal standard deviation of the accelerator
multiplier does not fall as $\lambda_u$ decreases across $[0.25, 1.0]$, and
reaches zero only at $\lambda_u = 0$. Utilization is highly persistent, and a
constant-gain filter removes the variance of white noise, not that of a
near-unit-root process, so smoothing an already-smooth signal changes little.
The turning point $\rhostar$, the whole-support slope and the wage-led reading
are correspondingly $\lambda_u$-invariant within inter-seed intervals across
$[0.25, 1.0]$, at every $\beta$ (\cref{tab:b17}). The finding is robust to how
the accelerator's signal is formed.

\begin{table}[htbp]
\centering
\small
\caption{Anchored-margin resolution across accelerator strength $\beta$ and
smoothing gain $\lambda_u$, at the anchored headline configuration
($c_0 = 1.0$, empirical $\sigma$, $\eta = 0.10$), with $\rho$ anchored at
$0.363$. Symbols: $\downarrow$ resolved on the falling branch (a negative
anchored margin); $\sim$ unresolved, the anchor lies inside the $\rhostar$
confidence interval (the margin vanishes); $\emptyset$ unresolved, the
$\rhostar$ interval is not estimable. Read from \texttt{ces\_b17\_margin};
\texttt{margin\_resolved} is authoritative, so an absent interval is never read
as resolved. The $\lambda_u = 0$ column switches the accelerator off: there
$\rhostar = 0.358$ for every $\beta$, with the anchor inside the interval.}
\label{tab:b17}
\begin{tabular}{lccccc}
\toprule
$\beta \,\backslash\, \lambda_u$ & $0.00$ & $0.25$ & $0.50$ & $0.75$ & $1.00$\\
\midrule  % generated by scripts/make_tab_b17.py -- do not hand-edit
$0.05$ & $\sim$ & $\sim$ & $\sim$ & $\sim$ & $\sim$\\
$0.15$ & $\sim$ & $\emptyset$ & $\emptyset$ & $\sim$ & $\sim$\\
$0.30$ & $\sim$ & $\downarrow$ & $\sim$ & $\sim$ & $\sim$\\
$0.50$ & $\sim$ & $\downarrow$ & $\downarrow$ & $\downarrow$ & $\downarrow$\\
$0.75$ & $\sim$ & $\downarrow$ & $\downarrow$ & $\downarrow$ & $\downarrow$\\
$1.00$ & $\sim$ & $\downarrow$ & $\downarrow$ & $\downarrow$ & $\downarrow$\\
\bottomrule
\end{tabular}
\end{table}

\paragraph{What this does and does not test.} Two boundaries hold whenever the
result is cited. First, the sweep is conditional on the anchored configuration:
it varies $\lambda_u$ and $\beta$ at that point, not marginally over the full
space, so it speaks to the anchored margin and not to the marginalized frontier
of \cref{sec:whatfails}; the marginal counterpart was not run. Second, it holds
$c_0 = 1.0$ and so never enters the $c_0 = 2.0$ collapse regime, where a
demand-expectation gain showed its only effect (\cref{sec:expectations}); that
omission is a design gap of this experiment, met instead by the price probe of
\cref{sec:wagecurve}, whose grid includes $c_0 = 2.0$.

\subsection{Firm heterogeneity: a cliff, not a gradient}
\label{sec:heterogeneity}

All firms share one productivity. That is an assumption, and rather than
implementing heterogeneity we measured what would happen if we did,
using a single experimental dial that fans firm productivities out
mean-preservingly, changing nothing else (no selection, no demand
reallocation, no entry or exit).

The result is a cliff (\cref{tab:heterogeneity}). Below a threshold no firm
dies; one grid step above it, every firm is dead, output is exactly
zero, unemployment is exactly one, and aggregate capital decays geometrically
to $3.5\times10^{-34}$ by step 2{,}000.

\begin{table}[htbp]
\centering
\small
\caption{Aggregates against productivity dispersion (as $U$\,/\,$Y$).}
\label{tab:heterogeneity}
\begin{tabular}{lcccr}
\toprule
Dispersion & S1 (anchor) & S2 (headline) & S3 ($rr = 0.5$) & Dead firms (S2)\\
\midrule
0.000 & 0.258 / 132.1 & 0.566 / 82.1 & 0.427 / 108.7 & 0\\
0.050 & 0.250 / 133.0 & 0.563 / 82.7 & 0.421 / 109.5 & 0\\
0.100 & 0.229 / 134.7 & 0.554 / 83.2 & 0.410 / 109.9 & 0\\
0.125 & 1.000 / 0.0 & 0.543 / 84.2 & 0.676 / 58.9 & 0\\
0.150 & 1.000 / 0.0 & 1.000 / 0.0 & 1.000 / 0.0 & 10\\
0.200 & 1.000 / 0.0 & 1.000 / 0.0 & 1.000 / 0.0 & 10\\
\bottomrule
\end{tabular}
\end{table}

\paragraph{The domino.} Tracing a collapsing run identifies the cascade: the
low-productivity firm serves the same network demand with more labor, earns
less profit, invests below $\delta K$ and decapitalizes first (its capital
falls from 38 to near zero by step 250). Its customers' spending shares remain
pointed at it, so that demand is destroyed rather than reallocated, and its laid-off workers lose their income, a demand externality. The
high-productivity firms follow: the strongest firm's capital reaches zero by
step 500. This is precisely what the missing machinery would have
prevented: with entry, exit and demand rerouting, that demand would have moved
to a surviving firm.

\paragraph{The benefit again fails as a cushion.} The hypothesis was that
keeping income flowing to the laid-off would raise the viability threshold. It
lowers it: at a dispersion of $0.125$ the headline scenario has 0 of 20
seeds with any dead firm, while the same scenario with $rr = 0.5$ has 18 of 20,
7 of them fully collapsed. The verified mechanism is the same channel once
more: the benefit lowers unemployment, the wage curve reads that and
raises the wage, and the low-productivity firm, whose marginal product
is scaled down by its own productivity, is the first pushed below
$I = \delta K$.

\paragraph{What this licenses us to claim.} Within-industry TFP dispersion is
large: \citet{Syverson2004} reports a 90/10 ratio near 2:1 in US manufacturing
(see also \citealp{BartelsmanDoms2000}). The collapse threshold here is a
max/min productivity ratio of about $1.22$--$1.29$. The units are not the
same and no quantitative mapping is claimed, a linear fan half-width is not
a 90/10 log-TFP ratio. The qualitative reading is all that is supported, and it
suffices: empirically realistic firm heterogeneity lies well outside this
model's viable range, because the model has no reallocation channel to absorb
it. Building heterogeneous firms without building reallocation would produce a
model that dies rather than a model with heterogeneous firms.

A collateral finding sharpens the mean-field question: even at zero dispersion
firms do not stay identical, because consumption links are random. The
capital share of the largest three firms starts at $0.30$ and settles at
$0.35$--$0.38$. The defensible statement is that the firm side is
quasi-representative in its aggregates, not in its cross-section.
\Cref{fig:heterogeneity} shows the cliff and the cascade together.

\begin{figure}[htbp]
\centering
\begin{subfigure}{0.49\textwidth}
  \includegraphics[width=\linewidth]{ces_b10_aggregates_spread.png}
  \caption{The cliff.}
\end{subfigure}\hfill
\begin{subfigure}{0.49\textwidth}
  \includegraphics[width=\linewidth]{ces_b10_domino_trace.png}
  \caption{The domino.}
\end{subfigure}
\caption{Firm heterogeneity: aggregates are stable until they are not.}
\label{fig:heterogeneity}
\end{figure}
